\documentclass[a4paper,11pt]{article}
\usepackage[utf8]{inputenc}
\usepackage[spanish]{babel}
\usepackage{geometry}
\usepackage{enumitem}
\usepackage{sectsty}
\usepackage{titlesec}
\usepackage{tikz}
\usepackage{array}

\geometry{
    a4paper,
    left=20mm,
    top=18mm,
    right=20mm,
    bottom=20mm
}

% Sin números de página
\pagestyle{empty}

% Estilo de secciones con números en círculos
\newcommand{\circled}[1]{%
    \tikz[baseline=(char.base)]{%
        \node[shape=circle,draw,inner sep=0.5pt,fill=black,text=white,font=\bfseries\footnotesize] (char) {#1};%
    }%
}

% Formato de secciones
\titleformat{\section}
{\normalfont\Large\bfseries}
{\circled{\thesection}}{0.5em}{}

\titlespacing*{\section}{0pt}{12pt}{8pt}

\titleformat{\subsection}
{\normalfont\normalsize\bfseries}
{}{0em}{}

\titlespacing*{\subsection}{8pt}{4pt}{4pt}

\begin{document}

% Encabezado del curso
\begin{center}
\small{Universidad de Granada}

\vspace{0.3cm}

\textbf{\large{Producción e Interacción Oral - Nivel 7}}

\vspace{0.4cm}

\small
\begin{tabular}{@{}ll@{}}
\textbf{Grupos:} & A (Aula 24) y B (Aula 08) \\
\textbf{Centro:} & Centro de Lenguas Modernas / Edificio A \\
\textbf{Profesor:} & Javier Benítez Lainez \\
\textbf{Tutorías:} & Martes 9:00-10:00 y 14:30-15:30 \\
& Sala de profesores / Email: benitezl@go.ugr.es
\end{tabular}
\end{center}

\vspace{0.5cm}
\section*{Conectores de Seguimiento}

\vspace{0.4cm}

\section{Para Mantener y Conectar la Conversación}

\vspace{0.2cm}

---

\vspace{0.2cm}

\section{¿QUÉ SON LOS CONECTORES DE SEGUIMIENTO?}

\vspace{0.2cm}

Los conectores de seguimiento son expresiones que permiten:

\begin{itemize}

  \item Reconocer lo que el interlocutor ha dicho

  \item Conectar tu intervención con lo anterior

  \item Seguir el hilo de la conversación

  \item Mostrar atención y escucha activa

  \item Tomar turnos apropiadamente

\end{itemize}

\vspace{0.2cm}

---

\vspace{0.2cm}

\section{CATEGORÍA 1: RECONOCIMIENTO}

\vspace{0.2cm}

\subsection{Para mostrar que estás escuchando}

\vspace{0.2cm}

\section{Confirmación básica}

\begin{itemize}

  \item \textbf{"Entiendo."} - Entiendo.

  \item \textbf{"Comprendo."} - Comprendo.

  \item \textbf{"Ya veo."} - Ya veo.

  \item \textbf{"Ah, entiendo."} - Ah, entiendo.

  \item \textbf{"Te entiendo."} - Te entiendo.

\end{itemize}

\vspace{0.2cm}

\section{Confirmación con matiz}

\begin{itemize}

  \item \textbf{"Entiendo lo que quieres decir."} - Entiendo lo que quieres decir.

  \item \textbf{"Veo por dónde vas."} - Veo por dónde vas.

  \item \textbf{"Entiendo tu punto."} - Entiendo tu punto.

  \item \textbf{"Comprendo tu posición."} - Comprendo tu posición.

  \item \textbf{"Entiendo perfectamente."} - Entiendo perfectamente.

\end{itemize}

\vspace{0.2cm}

\section{Validación}

\begin{itemize}

  \item \textbf{"Tienes razón en eso."} - Tienes razón en eso.

  \item \textbf{"Es un buen punto."} - Es un buen punto.

  \item \textbf{"Es muy interesante lo que dices."} - Es muy interesante lo que dices.

  \item \textbf{"Nunca había pensado en eso."} - Nunca había pensado en eso.

  \item \textbf{"Esa es una perspectiva interesante."} - Esa es una perspectiva interesante.

\end{itemize}

\vspace{0.2cm}

---

\vspace{0.2cm}

\section{CATEGORÍA 2: CONEXIÓN DIRECTA}

\vspace{0.2cm}

\subsection{Para conectar directamente con lo dicho}

\vspace{0.2cm}

\section{Parafrasear}

\begin{itemize}

  \item \textbf{"O sea que, según lo que dices..."} - O sea que, según lo que dices...

  \item \textbf{"Entonces, estás diciendo que..."} - Entonces, estás diciendo que...

  \item \textbf{"Si te entiendo bien,..."} - Si te entiendo bien,...

  \item \textbf{"Quieres decir que..."} - Quieres decir que...

  \item \textbf{"En otras palabras,..."} - En otras palabras,...

\end{itemize}

\vspace{0.2cm}

\section{Responder directamente}

\begin{itemize}

  \item \textbf{"En relación con lo que acabas de decir..."} - En relación con lo que acabas de decir...

  \item \textbf{"Respecto a tu punto sobre..."} - Respecto a tu punto sobre...

  \item \textbf{"Sobre lo que mencionabas antes..."} - Sobre lo que mencionabas antes...

  \item \textbf{"Siguiendo con tu idea de..."} - Siguiendo con tu idea de...

  \item \textbf{"Continuando con lo que decías..."} - Continuando con lo que decías...

\end{itemize}

\vspace{0.2cm}

---

\vspace{0.2cm}

\section{CATEGORÍA 3: AMPLIACIÓN}

\vspace{0.2cm}

\subsection{Para expandir sobre lo dicho}

\vspace{0.2cm}

\section{Pedir elaboración}

\begin{itemize}

  \item \textbf{"¿Puedes elaborar más sobre eso?"} - ¿Puedes elaborar más sobre eso?

  \item \textbf{"¿A qué te refieres exactamente?"} - ¿A qué te refieres exactamente?

  \item \textbf{"¿Podrías darme un ejemplo?"} - ¿Podrías darme un ejemplo?

  \item \textbf{"¿Podrías explicar eso más?"} - ¿Podrías explicar eso más?

\end{itemize}

\vspace{0.2cm}

\section{Agregar perspectiva}

\begin{itemize}

  \item \textbf{"Añadiría a eso que..."} - Añadiría a eso que...

  \item \textbf{"Para complementar lo que dices..."} - Para complementar lo que dices...

  \item \textbf{"Además de eso,..."} - Además de eso,...

  \item \textbf{"A eso, añadiría que..."} - A eso, añadiría que...

  \item \textbf{"En esa misma línea,..."} - En esa misma línea,...

\end{itemize}

\vspace{0.2cm}

---

\vspace{0.2cm}

\section{CATEGORÍA 4: TRANSICIÓN SUAVE}

\vspace{0.2cm}

\subsection{Para cambiar de tema o perspectiva}

\vspace{0.2cm}

\section{Cambiar suavemente}

\begin{itemize}

  \item \textbf{"Cambiando un poco de tema..."} - Cambiando un poco de tema...

  \item \textbf{"Pasando a otro aspecto..."} - Pasando a otro aspecto...

  \item \textbf{"Por otro lado..."} - Por otro lado...

  \item \textbf{"Desde otra perspectiva..."} - Desde otra perspectiva...

  \item \textbf{"Si hablamos de algo diferente..."} - Si hablamos de algo diferente...

\end{itemize}

\vspace{0.2cm}

\section{Retornar al tema}

\begin{itemize}

  \item \textbf{"Volviendo a lo que decíamos antes..."} - Volviendo a lo que decíamos antes...

  \item \textbf{"Retomando el tema inicial..."} - Retomando el tema inicial...

  \item \textbf{"Si volvemos al punto principal..."} - Si volvemos al punto principal...

  \item \textbf{"Para volver a lo anterior..."} - Para volver a lo anterior...

\end{itemize}

\vspace{0.2cm}

---

\vspace{0.2cm}

\section{CATEGORÍA 5: ACUERDO/DESA CUERDO}

\vspace{0.2cm}

\subsection{Para responder con acuerdo o desacuerdo}

\vspace{0.2cm}

\section{Acuerdo}

\begin{itemize}

  \item \textbf{"Estoy totalmente de acuerdo contigo."} - Estoy totalmente de acuerdo contigo.

  \item \textbf{"Coincido plenamente con eso."} - Coincido plenamente con eso.

  \item \textbf{"Comparto tu opinión sobre eso."} - Comparto tu opinión sobre eso.

  \item \textbf{"Exactamente, eso es lo que yo pienso."} - Exactamente, eso es lo que yo pienso.

  \item \textbf{"Estoy de la misma opinión."} - Estoy de la misma opinión.

\end{itemize}

\vspace{0.2cm}

\section{Acuerdo parcial}

\begin{itemize}

  \item \textbf{"En cierto modo, estoy de acuerdo."} - En cierto modo, estoy de acuerdo.

  \item \textbf{"Estoy de acuerdo en parte."} - Estoy de acuerdo en parte.

  \item \textbf{"Hasta cierto punto tienes razón."} - Hasta cierto punto tienes razón.

  \item \textbf{"Comparto parte de tu opinión."} - Comparto parte de tu opinión.

\end{itemize}

\vspace{0.2cm}

\section{Desacuerdo respetuoso}

\begin{itemize}

  \item \textbf{"Entiendo tu punto, pero veo las cosas de otra manera."} - Entiendo tu punto, pero veo las cosas de otra manera.

  \item \textbf{"Aunque entiendo tu posición, yo diría que..."} - Aunque entiendo tu posición, yo diría que...

  \item \textbf{"Tienes un punto interesante, pero..."} - Tienes un punto interesante, pero...

  \item \textbf{"Desde mi perspectiva, las cosas son diferentes."} - Desde mi perspectiva, las cosas son diferentes.

\end{itemize}

\vspace{0.2cm}

---

\vspace{0.2cm}

\section{CATEGORÍA 6: SEGUIMIENTO TEMPORAL}

\vspace{0.2cm}

\subsection{Para hacer referencia a tiempos anteriores}

\vspace{0.2cm}

\section{Algo mencionado antes}

\begin{itemize}

  \item \textbf{"Como mencionabas antes..."} - Como mencionabas antes...

  \item \textbf{"Sobre lo que hablábamos al principio..."} - Sobre lo que hablábamos al principio...

  \item \textbf{"Retomando lo que dijiste hace un momento..."} - Retomando lo que dijiste hace un momento...

  \item \textbf{"Siguiendo con tu idea anterior..."} - Siguiendo con tu idea anterior...

\end{itemize}

\vspace{0.2cm}

\section{Preparar el futuro}

\begin{itemize}

  \item \textbf{"Más adelante podremos hablar de..."} - Más adelante podremos hablar de...

  \item \textbf{"En un momento volveremos a..."} - En un momento volveremos a...

  \item \textbf{"Guardemos eso para después..."} - Guardemos eso para después...

\end{itemize}

\vspace{0.2cm}

---

\vspace{0.2cm}

\section{EJEMPLOS DE CONVERSACIONES}

\vspace{0.2cm}

\subsection{Ejemplo 1: Conversación académica}

\vspace{0.2cm}

\textbf{A:} "Creo que la tecnología está transformando la educación."

\vspace{0.2cm}

\textbf{B:} \textbf{Entiendo lo que quieres decir.} ¿Podrías darme un ejemplo específico?

\vspace{0.2cm}

\textbf{A:} "Sí, el uso de IA en la enseñanza personalizada."

\vspace{0.2cm}

\textbf{B:} \textbf{Ah, entiendo.} En esa misma línea, yo añadiría que también cambia el rol del profesor.

\vspace{0.2cm}

\textbf{A:} \textbf{Exactamente, eso es lo que yo pienso.} El profesor pasa a ser un facilitador.

\vspace{0.2cm}

\textbf{B:} \textbf{Respecto a tu punto sobre el rol del profesor,} creo que es necesaria más formación.

\vspace{0.2cm}

---

\vspace{0.2cm}

\subsection{Ejemplo 2: Conversación informal}

\vspace{0.2cm}

\textbf{A:} "Ayer vi una película increíble."

\vspace{0.2cm}

\textbf{B:} \textbf{¿De verdad? ¿Cuál?}

\vspace{0.2cm}

\textbf{A:} "El nuevo filme de Almodóvar."

\vspace{0.2cm}

\textbf{B:} \textbf{Ah, entiendo.} Yo también la vi. \textbf{Estoy totalmente de acuerdo contigo,} fue fenomenal.

\vspace{0.2cm}

\textbf{A:} \textbf{¿Cuál fue tu parte favorita?}

\vspace{0.2cm}

\textbf{B:} \textbf{Sobre lo que mencionabas antes,} creo que el final fue perfecto.

\vspace{0.2cm}

---

\vspace{0.2cm}

\subsection{Ejemplo 3: Debate}

\vspace{0.2cm}

\textbf{A:} "El gobierno debería aumentar el salario mínimo."

\vspace{0.2cm}

\textbf{B:} \textbf{Veo por dónde vas.} Sin embargo, eso podría causar inflación.

\vspace{0.2cm}

\textbf{A:} \textbf{Entiendo tu punto, pero veo las cosas de otra manera.} La evidencia de otros países muestra lo contrario.

\vspace{0.2cm}

\textbf{B:} \textbf{Tienes un punto interesante.} ¿Podrías elaborar más sobre eso?

\vspace{0.2cm}

\textbf{A:} \textbf{Claro. Para complementar lo que dices,} miraríamos el ejemplo de Alemania.

\vspace{0.2cm}

\textbf{B:} \textbf{Si te entiendo bien,} estás diciendo que el aumento salarial no causó inflación allí.

\vspace{0.2cm}

\textbf{A:} \textbf{Exactamente.}

\vspace{0.2cm}

---

\vspace{0.2cm}

\section{EJERCICIOS PRÁCTICOS}

\vspace{0.2cm}

\subsection{Ejercicio 1: Completación}

Completa la conversación con conectores de seguimiento apropiados:

\vspace{0.2cm}

\textbf{A:} "Creo que deberíamos viajar más."

\textbf{B:} ___________ . ¿A qué tipo de viaje te refieres?

\textbf{A:} "Viajes culturales."

\textbf{B:} ___________ . Yo añadiría que también son importantes los viajes de naturaleza.

\vspace{0.2cm}

\subsection{Ejercicio 2: Transformación}

Transforma las respuestas para usar conectores de seguimiento:

\vspace{0.2cm}

1. "No estoy de acuerdo." → _____________________

2. "¿Puedes explicar más?" → _____________________

3. "Estoy de acuerdo." → _____________________

\vspace{0.2cm}

\subsection{Ejercicio 3: Role-play}

Practica en parejas con el tema "El impacto de las redes sociales". Usa al menos 10 conectores de seguimiento diferentes.

\vspace{0.2cm}

---

\vspace{0.2cm}

\section{RÚBRICA DE EVALUACIÓN C1}

\vspace{0.2cm}

\subsection{Excelente}

\begin{itemize}

  \item Uso variado y natural de múltiples categorías

  \item Respuestas muestran escucha activa

  \item Transiciones suaves entre temas

  \item Conexiones lógicas con lo anterior

  \item Registro apropiado al contexto

\end{itemize}

\vspace{0.2cm}

\subsection{Bueno}

\begin{itemize}

  \item Uso adecuado de varias categorías

  \item Respuestas muestran atención

  \item Transiciones aceptables

  \item Conexiones presentes

  \item Registro generalmente apropiado

\end{itemize}

\vspace{0.2cm}

\subsection{Suficiente}

\begin{itemize}

  \item Uso limitado de conectores

  \item Algunas respuestas muestran atención

  \item Transiciones básicas

  \item Conexiones simples

  \item Registro variable

\end{itemize}

\vspace{0.2cm}

---

\vspace{0.2cm}

\textbf{Sesión 19: Estrategias de Interacción}

\textbf{Nivel C1 - Español Oral Avanzado}

\vspace{0.2cm}

\end{document}
